\begin{center}
\textsc{Gauss} an \textsc{Olbers}.\quad Zeven, 8.~Juli 1824.
\end{center}

\par\noindent\hspace*{\parindent}Ihre beiden letzten gütigen Briefe habe ich richtig erhalten .\ .\ .\ .\ .\ .

\par\noindent\hspace*{\parindent}Der erstere hat mich hinsichtlich aller meiner Messungen sehr nieder\-%
geschlagen. Da Sie das Dreieck Bremen-Brüttendorf-Wentel \textls[120]{wegen des
zu spitzen Winkels an Bremen}, verwerfen, so brechen Sie dadurch zugleich
den Stab über die, wie es scheint, einzig mögliche Art, auf der anderen Seite
um Bremen herum zu kommen; denn in dem Dreieck Bremen-Bottel-Stein\-%
berg wird der Winkel in Bremen noch viel spitzer sein. Sie setzen zwar mit Ihrer
gewohnten Güte hinzu, dass doch bei jenem Dreieck die Genauigkeit \textls[120]{des-}
\par\clearpage
\noindent{\small\textbf{ZUR HANNOVERSCHEN TRIANGULATION.}}\hfill371\par\medskip
\noindent\textls[120]{wegen} nicht bedeutend leiden würde, weil ich in meine Messungen eine so
große Schärfe lege. Allein dieser Grund, dessen Wahrheit ich jetzt auf sich
beruhen lassen will, kann mich durchaus im geringsten \textls[120]{nicht} beruhigen.
Nach meinem Grundsatz soll man immer, so genau man nur kann, beobachten;
der Grad der Genauigkeit in den Beobachtungen, gleichviel wie groß oder
wie klein er sein mag, bedingt immer wieder den Grad der Genauigkeit, die
man von den Resultaten fordern darf, und die Genauigkeit der Beobachtungen
kann nach meiner Meinung ein an sich schlechtes Dreieck durchaus nicht
gut machen; wenigstens wäre sonst überflüssig gewesen, die übrigen guten
Dreiecke mit derselben Schärfe zu messen. Höchstens kann dadurch dann
das ganze System wieder in Parallele mit anderen an sich viel schlechteren Mes\-%
sungen zurückkommen.

\par\noindent\hspace*{\parindent}Ich habe es bisher für ein bloßes Vorurteil gehalten, wenn man Dreiecke
mit sehr kleinen Winkeln der \textls[120]{Genauigkeit} für nachteilig hielt, insofern
die den spitzen Winkeln gegenüberliegenden Seiten keine Übergangsseiten
abgeben; ich habe solche kleinwinklige Dreiecke \textls[120]{bloß} deswegen für minder
gut gehalten, weil man damit auf einmal nicht viel weiter kommt, also mehr
Zeit und Kosten gebraucht, als wenn man auf einmal viel fortschreiten kann;
und auch dieser Grund fällt ganz weg, wenn die Aufsuchung und Instand\-%
setzung eines großen Dreiecks vielleicht doppelt so viel Zeit kostet, als die
Messung zweier Dreiecke zusammen, die eben dahin führen, und wovon das
eine einen sehr spitzen Winkel hat. Ungeachtet dessen habe ich nicht ganz
nach diesem Prinzip gehandelt, sondern ein Dreieck mit einem kleinen Winkel
nie eher adoptiert, als bis ich \textls[120]{fast alle} Möglichkeiten erschöpft hatte, es zu
vermeiden (nur diejenige Möglichkeit nicht, die zu schlechten Messungen
selbst geführt hätte, d.~i.\ \textsc{Zachs} hohe Türme); nicht weil ich geglaubt hätte;
dadurch an \textls[120]{Genauigkeit} etwas zu gewinnen, sondern aus dem wohl verzeih\-%
lichen Wunsche, dem System so viel möglich, außer dem inneren Gehalt, auch
Schönheit und Abrundung zu geben.

\par\noindent\hspace*{\parindent}Da ich nun aber \textls[120]{Sie} durch das, was ich in einem früheren Briefe
darüber schrieb, nicht überzeugt habe, sondern da Sie den Nachteil, der für
die Genauigkeit aus dem spitzen Winkel sonst entstehen würde, durch die
Schärfe der Messungen gut gemacht verlangen, was nach meiner Ansicht
\textls[120]{unmöglich} ist, so werde ich selbst in meiner bisherigen Ansicht ganz irre
\par\smallskip\noindent\hfill{\footnotesize 47\textsuperscript{*}}\par
\clearpage
\noindent372\hfill{\small\textbf{BRIEFWECHSEL MIT OLBERS.}}\par\medskip
\noindent und zweifelhaft, ob sie nicht ganz unrichtig gewesen ist, und darf wenigstens
auf keinen Fall hoffen, andere von der Richtigkeit derselben zu überzeugen. Was
namentlich das Dreieck Bremen-Brüttendorf-Wentel betrifft, so hätte ich mich
selbst sehr ungern dazu entschlossen, weil es nicht \textls[120]{schön} ist, und auf ein\-%
mal nicht viel weiter bringt; hinsichtlich der \textls[120]{Genauigkeit} aber (ganz
abgesehen davon, wie genau die Winkelmessungen an sich sind), würde ich
dasselbe, seine Winkel zu \(12^{0},46^{0},122^{0}\) angenommen, \textls[120]{vollkommen} einem
anderen gleichgestellt haben, dessen Winkel \(76^{0},46^{0},58^{0}\) gewesen wären.
